\begin{center}
\begin{tikzpicture}[x=.112\linewidth,y=7mm,every node/.style={font=\scriptsize,align=center}]
\fill[chFourTint] (0,3) rectangle (6.2,4);
\draw (0,0) rectangle (8.5,4);
\draw (1.5,0)--(1.5,4) (4.0,0)--(4.0,4) (6.2,0)--(6.2,4);
\draw (0,3)--(6.2,3) (1.5,2)--(4.0,2) (0,1)--(4.0,1);
\node at (.75,3.5) {Id no};
\node at (2.75,3.5) {Product name};
\node at (5.1,3.5) {Type};
\node at (.75,2) {1001};\node at (.75,.5) {1002};
\node at (2.75,2.5) {Computer};\node at (2.75,1.5) {Printer};\node at (2.75,.5) {Mouse};
\node at (5.1,1.5) {Accessories};
\node[draw,dashed,minimum width=10mm,minimum height=9mm] at (7.35,2) {ছবি};
\end{tikzpicture}
\captionof{figure}{অনুশীলন 20: পণ্যের সঙ্গে চার সারি জুড়ে ছবির ঘর}
\end{center}
\textbf{চিত্র ব্যাখ্যা:} ছবির ঘরটি শিরোনামসহ চার সারি জুড়ে আছে; এর জন্য আলাদা Image শিরোনাম লেখা হয়নি। 1001 ঘর দুই পণ্যের সারি এবং Accessories ঘর তিন পণ্যের সারি জুড়ে বিস্তৃত। ছকে ছবির স্থান বোঝানো হয়েছে; কোড চালালে \texttt{google.jpg} ফাইলের প্রকৃত ছবি সেখানে দেখা যাবে।
